\chapter{Reproduction and scientific code audit}
\section{Original implementation}
The supplied implementation puts input reading, replica separation, plotting, clustering and summary calculations in one script. Parsing the source identifies an Italian sentence without a comment marker at line 437, causing a syntax error before the analysis can run. The original is preserved exactly, including this defect, so that the audit can be checked against the supplied file.

The audit traces each calculation from its input to the quantity it would report. Reading the source shows what the script attempts to do, but no original execution manifest or complete numerical output set was located. Where the code differs from the paper, the method map records both versions and explains the correction made in the modular implementation.

\section{Identified methodological issues}
The original parser silently takes the first three segments separated by time resets and skips certain malformed lines. The new parser checks both the segment count and each row. Original group generation depends on a fixed list of 17 seeds, which would run out with smaller groups. The replacement uses local RNG streams and preserves every frame's identity. It makes the sampling reproducible but changes the exact subsets, so it is not claimed to reproduce the legacy script bit for bit.

The handling of residual frames changes which observations define a family. In the supplied code, rejected frames receive a generic label and every label enters consensus. Frames that failed the acceptance criteria can therefore influence the global centers. The modular version uses only accepted local clusters to discover families, then assigns the residual frames in STEP 3. Residuals still contribute to the final population without being treated as an accepted local cluster.

The original FES reader ignores the header information needed to check completeness and permits missing grid values with a warning. The weight function uses the nearest grid point and selects a boundary value for out-of-range queries. These choices can produce plausible-looking weights from invalid inputs. The replacement rejects incomplete grids and uses bilinear interpolation only within the grid bounds. The original partial download ended inside a numerical row and was rejected. The complete replacement received on 7 October passed grid validation; the earlier partial file remains unusable.

The original analysis weights cluster relevance but uses ordinary arithmetic means and standard deviations for the structural descriptors. Those summaries describe a different distribution from the weighted one specified in the paper's figure captions. The revised routines use the same weights for relevance and descriptor summaries within each cluster, and normalize replicas separately before pooling. Distribution width, replica SEM and bootstrap intervals remain separate quantities.

\begin{auditnote}
Correcting a calculation does not validate its input data or establish that the original simulation converged. The complete grid now passes validation, but its energy unit and construction procedure still need clarification before reweighting the real data.
\end{auditnote}

\section{Scientific V2 corrections: evidence unavailable}
The requested intermediate file \code{CLUSTERING40\_scientific\_v2.py} and its changelog were not located. The method map therefore marks that implementation as unavailable. None of the corrections in this report is attributed to the missing version. If supplied later, it should be checked by its hash, preserved unchanged, and compared with both the original and the modular implementation.

\section{Verification and numerical tests}
The current suite passed 36 tests, including strict parsing, interval boundaries, normalization, analytic bilinear interpolation, incomplete and duplicate grids, finite/bounds checks, weight sign, free-energy-offset invariance, equal replica masses, weighted moments, residual handling, final assignment, label-permutation agreement, and bootstrap separation of replicas. A seeded autoregressive fixture checks that the correlation estimator behaves plausibly for a known stationary process.

A complete synthetic test dataset exercises cluster discovery, consensus, all-frame assignment, final-bias weights, population tables, conditional bootstrap intervals and PDF generation. Its generated descriptors and analytic surface are used only for testing, not as Aβ40 observations. A further test checks repeatability by fitting the GMM extension to synthetic, separated observations. Tests with renamed labels and missing families check that the comparisons behave as intended. An archive-provenance test also checks that a source release cannot inherit the Git commit of an unrelated enclosing repository. These tests check numerical behavior; they cannot establish that a cluster represents a biological family.

\section{Remaining methodological ambiguities}
The main unresolved choices are the exact descriptor scaling in STEP 2, how the replica-averaged FES was constructed, and how the published cluster widths and uncertainties were defined. The supplied code also gives more detail than the paper about density ties and spline-based elbow selection. The repository identifies these choices as based on the supplied source; it does not imply that the paper fully specifies them.

These choices affect how the analysis is run and how its results can be compared. Raw Wasserstein scaling remains the reproduction setting; normalized scaling has a separate extension configuration. The complete numerical grid does not resolve how the FES was constructed. Any future difference in populations will need to be checked against these analysis choices before it is given a biological interpretation.
